Because (VaR, ES) is jointly elicitable, a neural network can be trained to forecast Expected Shortfall in exactly the way it is evaluated, by minimising a strictly consistent scoring function. We derived the properties that make this practical. Pooling across markets is exact under scale standardisation, and an output layer can guarantee coherent, non-crossing forecasts. We then evaluated the approach against a demanding set of benchmarks over twenty years and twelve markets.

Three conclusions follow. First, a pooled LSTM trained on FZ0 matches, but does not beat, the best econometric models (GJR-GARCH with skewed-$t$ innovations, filtered historical simulation and GARCH-EVT). Second, the deep models carry information that the econometric models do not. Adding them to an econometric forecast combination improves accuracy at every risk level, in almost every market and steadily over time. The result is the best model in our study and a modest saving in regulatory capital at an unchanged backtesting record. Third, the networks' weaknesses are specific and interpretable. They are slightly less conservative than the best econometric models and are slower to adapt to unprecedented volatility spikes. A hybrid that takes its volatility scale from GARCH mitigates the second weakness but learns little that improves forecasts. This suggests that, from daily returns alone, the predictable variation in tail shape beyond volatility is small.

\paragraph{Limitations.} The information set is deliberately restricted to daily returns, the horizon is one day, and the study covers equity indices only. Hyperparameters were fixed in advance rather than tuned. Two implementation choices were revised after a first run, both transparently documented and neither affecting the conclusions. The calibration test we use is mildly oversized in samples of the size of a single market's history.

\paragraph{Extensions.} Natural next steps include adding realised measures or intraday ranges to separate the value of information from that of architecture, multi-day horizons, and robustness of the ranking across the whole class of consistent scoring functions via Murphy diagrams \citep{ziegel2020robust}. The most interesting extension is economic: whether the time-varying tail shape learned by the hybrid is priced in the cross-section or predicts market returns, in the spirit of \citet{kelly2014tail}.
